% !TeX spellcheck = en_US
\documentclass[french]{yLectureNote}

\title{Mecanique Analytique}
\subtitle{Physique}
\author{Paulhenry Saux}
\date{\today}
\yLanguage{Français}
%lionels.calmels@cemes.fr
\professor{M.Brut}
\usepackage{graphicx}%----pour mettre des images
\usepackage[utf8]{inputenc}%---encodage
\usepackage{geometry}%---pour modifier les tailles et mettre a4paper
%\usepackage{awesomebox}%---pour les boites d'exercices, de pbq et de croquis ---d\'esactiv\'e pour les TP de PC
\usepackage{tikz}%---pour deiffner + d\'ependance de chemfig
% \usepackage{tabularx}%---pour dimensionner automatiquement les tableaux avec variable X
\usepackage{awesomebox}%---Pour les boites info, danger et autres
\usepackage{menukeys}%---Pour deiffner les touches de Calculatrice
\usepackage{fancyhdr}%---pour les en-t\^ete personnalis\'ees
\usepackage{blindtext}%---pour les liens
\usepackage{hyperref}%---pour les liens (\`a mettre en dernier)
\usepackage{caption}%---pour la francisation de la l\'egende table vers Tableau
\usepackage{pifont}
\usepackage{array}%---pour les tableaux
\usepackage{lipsum}
\usepackage{yFlatTable}
\usepackage{multicol}
\usepackage{tabularx}
\usepackage{booktabs}
\newcommand{\Lim}[1]{\lim\limits_{\substack{#1}}\:}
\renewcommand{\vec}{\overrightarrow}
\newcommand{\N}[0]{\mathbb{N}}
\newcommand{\dd}{\mathrm{d}}
\newcommand{\norm}[1]{||\vec{#1}||}
\newcommand{\fo}{\psi(\vec{r},t)}
\newcommand{\foe}{\psi(\vec{r},t)\*}
\newcommand{\HH}{\hat{H}}
\newcommand{\hb}{\hbar}
\newcommand{\lap}{\nabla^2}
\newcommand{\lapcc}{\frac{\partial^2 }{\partial x^2}+\frac{\partial^2 }{\partial y^2}+\frac{\partial^2 }{\partial z^2}}
\newcommand{\E}{\vec{E}(M)}
\newcommand{\B}{\vec{B}(M)}
\newcommand{\V}{V(M)}
\newcommand{\er}{\vec{e_{\rho}}}
\newcommand{\ep}{\vec{e_{\varphi}}}
\newcommand{\pip}{\pi^+}
\newcommand{\pim}{\pi^-}
\newcommand{\mv}{\mathcal{V}}
\DeclareMathOperator{\Div}{Div}
\newcommand{\rot}{\vec{\mathrm{rot}}}
\newcommand{\grad}{\vec{\mathrm{grad}}}
\newcommand{\qa}{q_{\alpha}}
\newcommand{\qdb}{\dot{q_{\beta}}}
\newcommand{\qb}{q_{\beta}}
%\setcounter{chapter}{3}
\begin{document}
\chapter{Mécanique Lagrangienne - Méthode}
\section{Méthode générale}
\subsection{Coordonnées généralisées}
\begin{definition}[Coordonnées généralisées]
Pour un système à $N$ corps possédant $k$ liaisons indépendantes, le nombre de degrés de liberté (ddl) est $n = 3N - k$. On choisit alors un ensemble de $n$ variables indépendantes $q = (q_1, q_2, \dots, q_n)$ appelées \textbf{coordonnées généralisées} permettant de décrire de manière unique la position de tout point du système.
\end{definition}


\checkInfo{Méthode de paramétrage}{
Une fois les $q_i$ choisis, exprimer systématiquement les coordonnées cartésiennes $(x_j, y_j, z_j)$ de chaque masse $m_j$ en fonction des $q_i$. 
Puis, dériver-les par rapport au temps pour obtenir les vitesses $(\dot{x}_j, \dot{y}_j, \dot{z}_j)$.
}
\subsection{Construction du Lagrangien}
\begin{proposition}[Énergie cinétique $T$]
L'énergie cinétique globale est la somme des énergies cinétiques de chaque masse. 
\begin{itemize}
    \item En cartésiennes : $T = \frac{1}{2}m(\dot{x}^2 + \dot{y}^2 + \dot{z}^2)$
    \item En polaires : $T = \frac{1}{2}m(\dot{r}^2 + r^2\dot{\theta}^2)$
    \item En sphériques : $T = \frac{1}{2}m\big(\dot{r}^2 + r^2\dot{\theta}^2 + r^2\sin^2\theta\,\dot{\phi}^2\big)$
\end{itemize}
\end{proposition}

\begin{proposition}[Énergie potentielle $V$]
Exprimer les forces conservatives qui s'appliquent. Les deux plus fréquentes sont :
\begin{itemize}
    \item La pesanteur : $V_{pes} = \pm mgz$ (le signe dépend de l'orientation de l'axe vertical).
    \item Un ressort de raideur $k$ et de longueur à vide $l_0$ : $V_{el} = \frac{1}{2}k(l - l_0)^2$.
\end{itemize}
\end{proposition}

\begin{definition}[Lagrangien]
Le Lagrangien du système est défini par :
$$L(q_i, \dot{q}_i, t) = T - V$$
\end{definition}
\subsection{Équations du mouvement}
\begin{theorem}[Équations d'Euler-Lagrange]
Pour chaque coordonnée généralisée $q_i$ (avec $i = 1, \dots, n$), l'équation du mouvement s'obtient par :
$$\frac{d}{dt}\left( \frac{\partial L}{\partial \dot{q}_i} \right) - \frac{\partial L}{\partial q_i} = 0$$
\end{theorem}

\checkInfo{Déroulé du calcul}{
Pour chaque variable $q_i$ :
\begin{enumerate}
    \item Calculer la quantité conjuguée (moment généralisé) $p_i = \frac{\partial L}{\partial \dot{q}_i}$.
    \item Dériver cette quantité par rapport au temps $\frac{d}{dt}(p_i)$ (attention aux dérivées composées !).
    \item Calculer la force généralisée $\frac{\partial L}{\partial q_i}$.
    \item Assembler le tout pour obtenir une équation différentielle du second degré.
\end{enumerate}
}
\subsection{Constantes du mouvement}
\subsubsection{Variables cycliques (Théorème de Noether)}

\begin{definition}[Variable cyclique]
Une coordonnée généralisée $q_i$ est dite \textbf{cyclique} si elle n'apparaît pas explicitement dans le Lagrangien, c'est-à-dire si :
$$\frac{\partial L}{\partial q_i} = 0$$
\end{definition}

\begin{proposition}[Conservation de l'impulsion]
Si $q_i$ est une variable cyclique, alors le moment généralisé conjugué $p_i$ est une \textbf{constante du mouvement} :
$$p_i = \frac{\partial L}{\partial \dot{q}_i} = \text{Constante}$$
\end{proposition}

\subsubsection{Conservation de l'énergie et Formule de Beltrami}

\begin{definition}[Intégrale de Jacobi / Énergie constante]
Si le Lagrangien ne dépend pas explicitement du temps ($\frac{\partial L}{\partial t} = 0$), alors la fonction énergie $h$ (ou intégrale de Jacobi) est conservée. Elle se calcule par :
$$h = \sum_{i=1}^n \dot{q}_i \frac{\partial L}{\partial \dot{q}_i} - L = \text{Constante}$$
Si les liaisons sont indépendantes du temps et que $V$ ne dépend pas des vitesses, alors $h = T + V = E_m$ (Énergie mécanique).
\end{definition}

\begin{theorem}[Formule de Beltrami]
Pour un système à \textbf{un seul degré de liberté} $q$ où $\frac{\partial L}{\partial t} = 0$, l'équation d'Euler-Lagrange s'intègre directement sous la forme :
$$L - \dot{q}\frac{\partial L}{\partial \dot{q}} = \text{Constante}$$
Cette formule est un cas particulier de la conservation de l'intégrale de Jacobi.
\end{theorem}
\section{Exemple - Machine d'Artwood}
La corde est inextensible et de longueur totale $L + \pi R$. La portion de corde enroulée sur la poulie est une demi-circonférence, soit $\pi R$. Les portions verticales mesurent $z_1$ et $z_2$. On en déduit la relation de contrainte holonome :
$$z_1 + z_2 + \pi R = L + \pi R \implies z_1 + z_2 = L$$
Il y a 2 coordonnées et 1 équation de contrainte, le système possède donc :
$$n = 2 - 1 = 1 \text{ degré de liberté}$$
On choisit $z_1$ comme coordonnée généralisée unique. On a alors $z_2 = L - z_1$, ce qui implique par dérivation temporelle : $\dot{z}_2 = -\dot{z}_1$.

\subsection{Construction du Lagrangien}
\textbf{Énergie cinétique ($T$) :} 
Les masses ont des mouvements purement verticaux:
$$T = \frac{1}{2}m_1\dot{z}_1^2 + \frac{1}{2}m_2\dot{z}_2^2 = \frac{1}{2}(m_1 + m_2)\dot{z}_1^2$$

\textbf{Énergie potentielle ($V$) :}
L'axe $(Oz)$ est vertical descendant, dans le sens de la gravité $\vec{g}$. L'énergie potentielle de pesanteur s'écrit donc avec un signe moins :
$$V = -m_1gz_1 - m_2gz_2 = -m_1gz_1 - m_2g(L - z_1) = -(m_1 - m_2)gz_1 - m_2gL$$

\textbf{Lagrangien ($L$) :}
$$L(z_1, \dot{z}_1) = T - V = \frac{1}{2}(m_1 + m_2)\dot{z}_1^2 + (m_1 - m_2)gz_1 + m_2gL$$
\subsection{Équation du mouvement}

On applique l'équation pour la variable $z_1$ :
$$\frac{d}{dt}\left(\frac{\partial L}{\partial \dot{z}_1}\right) - \frac{\partial L}{\partial z_1} = 0$$
\begin{itemize}
    \item $\frac{\partial L}{\partial \dot{z}_1} = (m_1 + m_2)\dot{z}_1 \implies \frac{d}{dt}\left(\frac{\partial L}{\partial \dot{z}_1}\right) = (m_1 + m_2)\ddot{z}_1$
    \item $\frac{\partial L}{\partial z_1} = (m_1 - m_2)g$
\end{itemize}
On obtient l'équation différentielle du mouvement :
$$(m_1 + m_2)\ddot{z}_1 - (m_1 - m_2)g = 0 \implies \ddot{z}_1 = \frac{m_1 - m_2}{m_1 + m_2}g$$
\section{Double tige}
Deux masselottes considérées comme ponctuelles se déplacent sans frottements sur une table que l’on identifie au plan (O, x, y). La première masselotte,
de masse $m_1$, est reliée à l’origine par un première tige rigide de masse négligeable
et de longueur $l_1$. La deuxième masselotte, de masse $m_2$, est attachée à la première par une deuxième tige rigide elle aussi de masse négligeable et de longueur
$l_2$.

On appelle $\theta$ l’angle que fait la première tige fait avec l’axe $O_x$ et $\alpha$ l’angle
que fait la deuxième tige avec la première.

Les longueurs des tiges $l_1$ et $l_2$ sont rigides et constantes. Les seules variables indépendantes sont les angles $\theta$ et $\alpha$. Le système possède donc :
$$n = 2 \text{ degrés de liberté}$$
Les coordonnées généralisées choisies sont $q = (\theta, \alpha)$.


L'angle absolu de la deuxième tige par rapport à l'axe $(Ox)$ est $(\theta + \alpha)$. Par projection trigonométrique :
\begin{itemize}
    \item \textbf{Pour $m_1$ :}
    $$x_1 = l_1\cos\theta, \quad y_1 = l_1\sin\theta$$
    \item \textbf{Pour $m_2$ :}
    $$x_2 = l_1\cos\theta + l_2\cos(\theta+\alpha), \quad y_2 = l_1\sin\theta + l_2\sin(\theta+\alpha)$$
\end{itemize}
\subsection{Construction du Lagrangien}
\subsubsection{Calcul de T}
Pour calculer $T$, on dérive d'abord les positions par rapport au temps :
\begin{itemize}
    \item $\dot{x}_1 = -l_1\dot{\theta}\sin\theta \quad \text{et} \quad \dot{y}_1 = l_1\dot{\theta}\cos\theta \implies \dot{x}_1^2 + \dot{y}_1^2 = l_1^2\dot{\theta}^2$
    \item $\dot{x}_2 = -l_1\dot{\theta}\sin\theta - l_2(\dot{\theta}+\dot{\alpha})\sin(\theta+\alpha)$
    \item $\dot{y}_2 = l_1\dot{\theta}\cos\theta + l_2(\dot{\theta}+\dot{\alpha})\cos(\theta+\alpha)$
\end{itemize}
En développant $\dot{x}_2^2 + \dot{y}_2^2$ et en utilisant l'identité $\cos a\cos b + \sin a\sin b = \cos(a-b)$, on obtient :
$$\dot{x}_2^2 + \dot{y}_2^2 = l_1^2\dot{\theta}^2 + l_2^2(\dot{\theta}+\dot{\alpha})^2 + 2l_1l_2\dot{\theta}(\dot{\theta}+\dot{\alpha})\cos\alpha$$


L'énergie cinétique totale $T = T_1 + T_2$ s'écrit :
$$T = \frac{1}{2}(m_1+m_2)l_1^2\dot{\theta}^2 + \frac{1}{2}m_2l_2^2(\dot{\theta}+\dot{\alpha})^2 + m_2l_1l_2\dot{\theta}(\dot{\theta}+\dot{\alpha})\cos\alpha$$
\subsubsection{Lagrangien}

Le mouvement a lieu sur une table horizontale plane $(O,x,y)$. En l'absence de forces verticales ou de gravité mentionnée dans ce plan, l'énergie potentielle est constante : $V = 0$. On a donc :
$$L(\theta, \alpha, \dot{\theta}, \dot{\alpha}) = T$$
\subsection{Variable cyclique et quantités conservées}

En observant l'expression de $L$, on constate que l'angle $\theta$ n'apparaît nulle part explicitement (seule sa dérivée $\dot{\theta}$ est présente).
$$\frac{\partial L}{\partial \theta} = 0 \implies \theta \text{ est une variable cyclique.}$$


Puisque $\theta$ est cyclique, le moment généralisé conjugué $p_\theta$ associé à cette coordonnée est une constante du mouvement :
$$p_\theta = \frac{\partial L}{\partial \dot{\theta}} = \text{Constante}$$
Calculons explicitement cette quantité :
$$p_\theta = (m_1+m_2)l_1^2\dot{\theta} + m_2l_2^2(\dot{\theta}+\dot{\alpha}) + m_2l_1l_2(2\dot{\theta}+\dot{\alpha})\cos\alpha = \text{Constante}$$
Cette constante représente le moment cinétique total du système par rapport à l'axe vertical passant par l'origine $O$.\marginInfo{On pouvait s'en douter car le système est invariant par rotation globale d'un angle $\theta$.} 

\end{document}

